\subsubsection{Segment Tree iterativo}
\begin{lstlisting}[style=mystyle, language=C++]
// TC: O(log N) per operation
// usa: cuando quieres evitar recursion (mas rapido en practice)
#include <bits/stdc++.h>
using namespace std;

struct SegTreeIter{
	int n;
	vector<long long> t;
	SegTreeIter(int n){
		this->n = n;
		t.assign(2*n, 0);
	}
	void build(const vector<long long>& a){
		for(int i=0;i<(int)a.size();i++) t[n+i] = a[i];
		for(int i=n-1;i>=1;i--) t[i] = t[2*i] + t[2*i+1];
	}
	void set(int p, long long val){
		p += n;
		t[p] = val;
		for(p >>= 1; p >= 1; p >>= 1) t[p] = t[2*p] + t[2*p+1];
	}
	long long query(int l, int r){
		long long res = 0;
		for(l += n, r += n+1; l < r; l >>= 1, r >>= 1){
			if(l & 1) res += t[l++];
			if(r & 1) res += t[--r];
		}
		return res;
	}
};

int main(){
	SegTreeIter st(5);
	st.build({1, 2, 3, 4, 5});
	cout << st.query(1, 3) << "\n";  // 9
	st.set(2, 10);
	cout << st.query(1, 3) << "\n";  // 16
	return 0;
}
\end{lstlisting}
